\section{コンパクト性}
何事も有限性のあるほうが扱いやすいという意味で, コンパクト性は大きな役割を果たす。なお, ここでいう「コンパクト」はHausdorff性を意味しない。

\begin{eg} (\tf{いろんなコンパクト空間})
\ben
\item $\mb{R}^N$の有界閉集合はコンパクトである(\ref{theo:heineborel})。たとえば$S^{N-1}$や$D^{N}$, $D^1 = [-1,1]$など。
\item 有限集合はコンパクトである。
\item \tf{コンパクト空間の有限の和集合はコンパクトである。}
\item \tf{コンパクト空間なら閉集合$\implies$コンパクト。}これは「Hausdorff空間ならコンパクト集合$\implies$閉集合」と混同しやすいかもしれないが, 後者の場合は$\mb{R}^{N}$のような非コンパクトHausdorff空間を考える。
\item
\een
\end{eg}

\begin{eg} (\tf{コンパクト空間ではないもの})
\ben
\item $\mb{R}^{n}, S^{n} - \{ p \}$, $(0,1)$
\item \tf{コンパクト空間の共通部分はコンパクトとは限らない。}
\item 多様体から$1$点を抜いたものを$M$とする。$M$はコンパクトではない。なぜなら, その1点のまわりでは局所的に$\mb{R}^{n}$と同相だが, その点を含む閉単位球の無限減少列(半径が$0$に収束)が取れる。その補集合は$M$の開集合の上昇列を与え, $M$を被覆する。しかし, そこから有限被覆を取ることはできない。
\een
\end{eg}

\begin{theo}

\end{theo}

\begin{theo} (\tf{コンパクト空間まとめ})\\
\ben
\item 距離空間において, コンパクトであることと点列コンパクトであることは同値である。
\item
\een
\end{theo}

次にコンパクト空間とその連続写像に対する性質をまとめる.

\lefba{
\begin{prop} \label{prop:compact_to_haus}
\ben
\item コンパクト空間からハウスドルフ空間への連続な全単射写像は同相写像である\footnote{この結果は覚えておくとよい。多様体論でも使われることがある命題でもあるようである。}。
\item コンパクト空間の連続写像による像はまたコンパクトである.
\item (\tf{最大値の定理}) コンパクト空間上の連続関数は最大値を持つ.
\item
\een
\end{prop}
}
\prf
(1): 定義をいろいろ振り返る良い演習問題になるので自分でやってみよう。
(2): $K\subset X$をコンパクト集合とせよ. $f(K)$の開被覆$\set{U_i}$を取り, これを$f$で引き戻すと$f^{-1}(U_i)$たちが$K$を被覆する. コンパクト性からこの逆像被覆から有限被覆が取れるので, そこに現れる添え字の$U_i$たちが$f(K)$をすでに被覆している.
(3):

\prf

\begin{eg}
コンパクトの順像がコンパクトはいろいろ便利である. たとえば, $X=\mb{R}/\mb{Z}$という位相空間は, 自然な写像$\pi:\mb{R}\to \mb{R}/\mb{Z}$について, $\pi([0,1]) = \mb{R}/\mb{Z}$を満たすので, $\mb{R}/\mb{Z}$はコンパクトである.
\end{eg}
\begin{prac}
次の空間がコンパクトであることを証明せよ.
\ben
\item $\mb{T}^2$
\item $\mb{CP}^n$
\een
\end{prac}

\begin{egprob}
コンパクト部分空間の共通部分はまたコンパクトになるか?
\end{egprob}
\begin{ans}
$X:=\mb{R}\cup\set{\pm \infty}$に,
\[\varnothing, X, \mb{R}の任意の部分集合, \mb{R}\cup\set{\infty}, \mb{R}\cup\set{-\infty}\]
によって位相が定まる.このとき$U,V$はコンパクトだが$\mb{R} = U\cap V$はコンパクトではない.
\end{ans}




\clearpage
\subsection{例題}
\begin{egprob} (2018集合と位相演習)\\
Hausdorff空間$X$の空でないコンパクト集合\footnote{$X$の部分集合で, 任意の開被覆が有限被覆を持つもの。設問的に誘導位相でコンパクトであることは認めてはならない。}族$\{ K_n \}_{n\in \mb{N}}$が任意の$n\in \mb{N}$で$K_n\supset K_{n+1}$を満たすとする。
\ben
\item $K_1$に誘導位相を与えるとき, $K_1-K_n$は$K_1$の開集合であることを示せ
\item $K_1$に誘導位相を与えるとき, $K_1$はコンパクトであることを示せ。
\item $\disp\bigcap_{n\in \mb{N}}K_n \neq \varnothing$を示せ。
\een
\end{egprob}
\begin{egans}
Hausdorff空間のコンパクト集合は閉集合であるから, $K_n$はすべて閉集合であることに注意。

(1): $K_1 - K_n = K_1 \cap (X-K_n)$で, $X-K_n$は開集合だから誘導位相の定義より従う。

(2): $K_1$の被覆$\bigcup K_1\cap U_{\lambda}$を取る$(U_{\lambda} \opensub X)$。$\bigcup U_{\lambda} \supset K_1$だから有限被覆が取れ, ただちに従う。

(3): 空だと仮定すると, $U_n = K_1 - K_n$ ($n=1,2,\cdots$)は$K_1$を被覆するが, $K_n\neq \varnothing$なのでこの$\{ U_n \}_{n=1}^{\infty}$から有限被覆を取ることはできず, (2)に反する。
\end{egans}





\begin{egprob} (少々難しい例)\\
$\mb{Z}$の素数$p$に対して,$d(x,y) = c^{-v_{p}(x-y)}$とおくと$\mb{Z}$は距離空間となる。この距離を$p$進距離という。距離空間であるから, $\mb{Z}$の完備化があり, それを$\mb{Z}_{p}$と書く。$p$進距離で$\mb{Z}_{p}$はコンパクトであることを示せ。
\end{egprob}
\begin{egans}
距離空間であるから,点列コンパクト性, すなわち任意の$\{ x_n \}_{n=1}^{\infty} \subset \mb{Z}_{p}$が収束部分列をもつことを示せばよい。なお, 自然な写像によって同型$\mb{Z}_{p}/p\mb{Z}_{p} \cong \mb{Z}/p\mb{Z}$が存在することは認めることにする(雪江整数論1命題9.1.31を参照されたい)。また, $R$を0を含む$\mb{Z}/p\mb{Z}$の完全代表系とする。

$\{ x_n \}_{n=1}^{\infty}\subset \mb{Z}_{p}$とする。認めた同型によって$\mb{Z}_{p}/p\mb{Z}_{p}$は有限なので, $x_n - a_0 \in p\mb{Z}_{p}$であるような$n$が無数に存在する。それらを集めて$\{ x_{1,n}\}$という部分列を取る。$x_{1,n} = a_0 + p y_{1,n}$をみたすような$y_{1,n}\in \mb{Z}_{p}$がある。同様に無限個の$n$に対して$y_{1,n} - a_1 \in p\mb{Z}_{p}$である。このような操作を繰り返して, $\{ x_{m-1,n}\}_{n=1}^{\infty}$の部分列$\{ x_{m,n}\}_{n=1}^{\infty}$ですべての$n$に対し
\[x_{m,n} \equiv a_0 + p a_1 + \cdots + p^{m-1}a_{m-1} \pmod{p^{m} }\]
となるものを取ることができる。あとは対角線論法である。$x_{n+1,n+1} = x_{n,l}$となるような$l$は$n+1$以上だから, $\{ x_{n,n}\}_{n=1}^{\infty}$は$\{ x_n \}$の部分列(ここで気にしているのは「逆流」が起こらないかどうかということである)であり
\[x_{n,n} \equiv a_0 + p a_1 + \cdots p^{N} a_N \pmod{p^{N+1}\mb{Z}_{p}}\quad (n>N)\]
なので
\[\disp\lim_{N\to \infty} x_{n,n} = a_{0} + p a_1 + \cdots p^{N}a_N \in \mb{Z}_{p} \]
であるから$\{ x_n \}$は$\mb{Z}_{p}$の点に収束する部分列を持つ。したがって$\mb{Z}_{p}$は$p$進距離でコンパクトである。

なお, $\mb{Z}_{p}$以外にもDedekind環(\ref{defn:dedekind_domain})と呼ばれる局所環で素イデアル$\maf{p}$を使って$\maf{p}$進完備化$\wiha{A}_{\maf{p}}$を構成する方法があり, 適当な「有限性の仮定」を設けることで$\wiha{A}_{\maf{p}}$がコンパクトであることの証明も同様にできる。興味がある方は\cite{momoyukie}を読もう。
\end{egans}

\clearpage
